\section*{الفصل \ref{ch:b2:conjugate-additions} --- الكربونيلات المترافقة: مايكل وروبنسون}

\begin{solution}{exo:b2:conjugate-additions:1}
\begin{itemize}
  \item \ce{CH3MgBr}: \omterm{def:b2:conjugate-additions:addition-modes}{إضافة 1،2} أساسًا (قاسٍ، غير عكوس)، 1-ميثيل هكس-2-ين حلقي-1-ول.
  \item \ce{(CH3)2CuLi}: إضافة 1،4، 3-ميثيل هكسانون حلقي.
  \item الثيوفينول مع قاعدة: إضافة 1،4 (الثيولات لين)، 3-(فينيل سلفانيل)هكسانون حلقي.
  \item \ce{LiAlH4}: \omterm{def:b2:conjugate-additions:addition-modes}{إضافة 1،2}، هكس-2-ين حلقي-1-ول.
\end{itemize}
\end{solution}

\begin{solution}{exo:b2:conjugate-additions:2}
المانحات: البنتان-2،4-ديون، والنتروميثان، وبروبانديوات ثنائي الإيثيل (رابطة \ce{C-H} حمضية). والمستقبلات:
البروبيننتريل، وبروبينوات الميثيل، والبوت-3-ين-2-ون (رابطة \ce{C=C} مترافقة مع مجموعة ساحبة).
\end{solution}

\begin{solution}{exo:b2:conjugate-additions:3}
\ce{C4H9Br + 2Li -> C4H9Li + LiBr}، ثم \ce{2C4H9Li + CuI -> (C4H9)2CuLi + LiI}، في الإيثر في درجة حرارة
منخفضة تحت غاز خامل.
\end{solution}

\begin{solution}{exo:b2:conjugate-additions:4}
(3-أوكسوبيوتيل)بروبانديوات ثنائي الإيثيل، \ce{(EtOOC)2CHCH2CH2COCH3}. وتعطي الحلمهة حمض
البروبانديويك المستبدَل، الذي يفقد \ce{CO2} بالتسخين (\cref{prop:b2:enolates-aldol:decarboxylation}):
حمض 5-أوكسو هكسانويك، \ce{HOOCCH2CH2CH2COCH3}.
\end{solution}

\begin{solution}{exo:b2:conjugate-additions:5}
ثنائي ميثيل كوبرات الليثيوم في الإيثر في درجة حرارة منخفضة، ثم معالجة مائية: \omterm{def:b2:conjugate-additions:addition-modes}{إضافة مترافقة}.
أما بروميد ميثيل المغنيزيوم، وهو أقسى، فيُضاف في معظمه \omterm{def:b2:conjugate-additions:addition-modes}{إضافة 1،2} ويعطي 1-ميثيل هكس-2-ين حلقي-1-ول ناتجًا
رئيسيًا.
\end{solution}

\begin{solution}{exo:b2:conjugate-additions:6}
ينزع \ce{Et3N} بروتونًا من قليل من \ce{EtSH} فيعطي \ce{EtS-}؛ ويُضاف الثيولات إلى الكربون $\beta$، 
ويأخذ \omterm{def:b2:enolates-aldol:enolate}{أيون الإينولات} المتكوّن بروتون جزيء آخر من \ce{EtSH}، مما يجدد \ce{EtS-}. والناتج:
4-(إيثيل سلفانيل)بيوتان-2-ون. والكبريت كبير وقابل للاستقطاب، وزوجه الحر مرتفع الطاقة: فهو \omterm{def:b2:frontier-orbitals:hard-soft}{محب
للنواة لين}، تحت \omterm{def:b2:frontier-orbitals:control}{التحكم المداري}، فيهاجم حيث يكون المدار \omterm{def:b2:frontier-orbitals:frontier}{LUMO} أكبر، والإضافة
عكوسة، مما يحابي أيضًا ناتج 1،4 الأكثر استقرارًا.
\end{solution}

\begin{solution}{exo:b2:conjugate-additions:7}
يهاجم السيانيد كربون الكربونيل الأكثر إيجابية أسرع: فالسيانوهيدرين هو الناتج الحركي. وتكوّنه
عكوس؛ ففي درجة الحرارة الأعلى، ومع الوقت، يتحرر السيانيد ويُضاف من جديد عند الكربون
$\beta$، معطيًا الناتج الذي يحتفظ برابطة $\pi$ الأقوى في \ce{C=O}، وهو الناتج الترموديناميكي.
\end{solution}

\begin{solution}{exo:b2:conjugate-additions:8}
مع 2-ميثيل هكسان حلقي-1،3-ديون: كيتون فيلاند--ميشر (شكل \omterm{def:b2:conjugate-additions:robinson}{الإلحاق الحلقي لروبنسون}).
ومع الهكسانون الحلقي: 4،4a،5،6،7،8-سداسي هيدرونفثالين-2(3\textit{H})-ون، وهو إينون ثنائي الحلقة تنتهي
رابطته \ce{C=C} عند كربون التحام الحلقتين، بلا مجموعة ميثيل زاوية.
\end{solution}

\begin{solution}{exo:b2:conjugate-additions:9}
\ce{CH3COCH2CH2CH2COCH3}: C2 وC6 هما الكربونيلان، وبينهما خمس ذرات كربون من C2 إلى C6 مع احتسابهما،
فهو ثنائي كيتون-1،5. اقطع C3--C4: المانح أيون إينولات البروبانون عند C3 (والأفضل 3-أوكسوبيوتانوات الإيثيل، الذي
تُزال مجموعة الإستر فيه لاحقًا بالحلمهة \omterm{def:b2:enolates-aldol:decarboxylation}{ونزع الكربوكسيل})؛ والمستقبل البوت-3-ين-2-ون. والظروف:
3-أوكسوبيوتانوات الإيثيل، والبوت-3-ين-2-ون، وإيثوكسيد الصوديوم بكمية حفزية في الإيثانول؛ ثم حمض مائي 
وتسخين.
\end{solution}

\begin{solution}{exo:b2:conjugate-additions:10}
يكوّن \omterm{def:b2:enolates-aldol:aldol}{الألدول} رابطة \ce{C-C} لكنه يحوّل رابطة $\pi$ في \ce{C=O} إلى رابطة $\sigma$ في \ce{C-O}؛ 
والحصيلة صغيرة وجزيئان يصيران جزيئًا واحدًا، فيكون $K$ متواضعًا \omterm{def:b2:enolates-aldol:aldol}{والألدول} العكسي سهلًا. أما
\omterm{def:b2:conjugate-additions:michael}{إضافة مايكل} فتكوّن رابطة \ce{C-C} على حساب رابطة $\pi$ أضعف في \ce{C=C} وتحتفظ بكلتا
رابطتي \ce{C=O}: فهي مواتية بوضوح. ويتطلب التفاعل العكسي أن يطرد أيون إينولات الكيتون أنيون البروبانديوات المثبت؛ 
ويقع التوازن بعيدًا جدًا في جهة الناتج بحيث يبقى الناتج قائمًا مع قاعدة حفزية.
\end{solution}

\begin{solution}{exo:b2:conjugate-additions:11}
\ce{C(COOEt)2(CH2CH2COOCH3)2}. بعد \omterm{def:b2:conjugate-additions:michael}{إضافة مايكل} الأولى يبقى للكربون المركزي
هيدروجين واحد بين مجموعتي إستر، حمضي كما كان: فتنزعه القاعدة وتتبع ذلك إضافة ثانية.
\end{solution}

\begin{solution}{exo:b2:conjugate-additions:12}
عبر \omterm{def:b2:enolates-aldol:enolate}{أيون الإينولات} الأكثر استبدالًا (نحو الكربون الحامل للميثيل): الأوكتالون الذي تقع فيه مجموعة الميثيل
على كربون الالتحام، 4a-ميثيل-4،4a،5،6،7،8-سداسي هيدرونفثالين-2(3\textit{H})-ون. وعبر
\omterm{def:b2:enolates-aldol:enolate}{أيون الإينولات} الأقل استبدالًا (من جهة \ce{CH2}): تنتهي مجموعة الميثيل على كربون حلقي مجاور
للالتحام. ويُفضَّل الأول في الظروف التي تسمح بالتوازن (ألكوكسيد في كحوله، أو طريق
\omterm{def:b2:amines:imine}{الإينامين} \omterm{def:b2:amines:class}{بأمين ثانوي} وحمض)، التي تعطي أيون \omterm{def:b2:enolates-aldol:kinetic-enolate}{الإينولات الترموديناميكي} الأكثر استبدالًا
(\cref{prop:b2:enolates-aldol:enolate-control}).
\end{solution}

\begin{solution}{pb:b2:conjugate-additions:1}
\textbf{1.} حلقة هكسان حلقي فيها \ce{C=O} عند C1 وC3 ومجموعة ميثيل على C2؛ والهيدروجين على C2
هو الأكثر حموضة.
\textbf{2.} نزعه يعطي أنيونًا تتوزع شحنته على C2 والأكسجينين كليهما؛ أما في
الهكسانون الحلقي فلا يتقاسمها إلا أكسجين واحد، والحموضة أضعف من أن تُقاس في الماء.
\textbf{3.} الشحنة على C2؛ و\ce{C1=C2} مع \ce{O^-} على C1؛ و\ce{C2=C3} مع \ce{O^-} على C3.
\textbf{4.} نعم: مع الهيدروكسيد، $K = 10^{14.0 - 5.3} = 10^{8.7}$ (الماء بوصفه الحمض المرافق، من
$K_e$)؛ ومع الإيثوكسيد، $10^{15.9-5.3} = 10^{10.6}$. وكلاهما تام.
% ledger: ke:25C, pka:ethanol
\textbf{5.} \ce{C1(OH)=C2(CH3)}: ما زال C2 يحمل هيدروجينًا واحدًا، وهو كل ما يحتاج إليه \omterm{def:b2:enolates-aldol:tautomerism}{الإينول}؛ ومجموعة
الميثيل لا تحل إلا محل الهيدروجين الثاني.
\textbf{6.} لين: فالشحنة موزعة على ثلاث ذرات والأنيون مثبت. ويهاجم
الكربون $\beta$ (\cref{prop:b2:conjugate-additions:regio}).
\textbf{7.} يرتبط C2 في \omterm{def:b2:enolates-aldol:enolate}{أيون الإينولات} بطرف \ce{CH2} في البوت-3-ين-2-ون؛ ويأخذ \omterm{def:b2:enolates-aldol:enolate}{أيون الإينولات} المتكوّن
بروتونًا: 2-ميثيل-2-(3-أوكسوبيوتيل)هكسان حلقي-1،3-ديون.
\textbf{8.} كل أيون إينولات جديد يأخذ بروتونًا من جزيء ديون (أو من المذيب) فيجدد بذلك
أيون إينولات المانح: فلا تُستهلك القاعدة.
\textbf{9.} \omterm{def:b2:enolates-aldol:enolate}{أيون الإينولات} ثنائي الموقع؛ ومع محب للإلكترونات كربوني لين، تحت \omterm{def:b2:frontier-orbitals:control}{التحكم المداري}، يتفاعل
عند الكربون، وناتج \ce{C-C} يحتفظ برابطتي \ce{C=O}، وهي التوليفة الأقوى؛ وناتج
\ce{O}، إن تكوّن، ينعكس.
\textbf{10.} واحدة: C2، المرتبطة بكل من C1 وC3 والميثيل والسلسلة.
\textbf{11.} كربون كربونيل السلسلة وكربونيل الحلقة C1 (أو C3): \babelsublr{\ce{C(=O)}--\ce{CH2}--\ce{CH2}--C2--C1}،
خمس ذرات كربون مع احتساب الطرفين.
\textbf{12.} لاستهلاك كل الديون، وهو الكاشف الأثمن؛ وطفيف فقط، لأن البوت-3-ين-2-ون
شديد السمية ويتبلمر.
\textbf{13.} C4 وC6 في الحلقة (المجاوران للكربونين C3 وC1)، و\ce{CH2} في السلسلة المجاور لكربونيلها، 
و\ce{CH3} الطرفي في السلسلة. أما C2 فلا هيدروجين له.
\textbf{14.} أيون إينولات \ce{CH3} الطرفي يهاجم كربونيلًا حلقيًا: فتضم الحلقة المغلقة ذلك
الكربون، وكربون كربونيل السلسلة، و\ce{CH2} الاثنين، وC2 وكربون الكربونيل المهاجَم: ست ذرات.
\textbf{15.} أيون إينولات \ce{CH2} الداخلي في السلسلة كان سيغلق حلقة رباعية متوترة. وأيون إينولات
حلقي (C4 أو C6) يهاجم كربونيل السلسلة يغلق حلقة سداسية أيضًا، لكنها جسرية: وذلك \omterm{def:b2:enolates-aldol:aldol}{الألدول}
لا يستطيع نزع الماء (إذ ستقع الرابطة المزدوجة عند رأس جسر نظام ثنائي الحلقة صغير) فينعكس،
بينما يُستنزف \omterm{def:b2:enolates-aldol:aldol}{الألدول} الملتحم بنزع الماء.
\textbf{16.} يحمل \omterm{def:b2:enolates-aldol:aldol}{الألدول} \ce{OH} على كربون كربونيل الحلقة السابق، الذي صار الآن عند التحام الحلقتين؛ وفقد
الماء مع هيدروجين من كربون \ce{CH3} السابق (الذي صار في الموضع $\alpha$ بالنسبة إلى كيتون السلسلة) يعطي
\ce{C=C} مترافقة مع ذلك الكيتون.
\textbf{17.} الرابطة المزدوجة الجديدة مترافقة مع الكربونيل: حذف E1cB عبر
\omterm{def:b2:enolates-aldol:enolate}{أيون الإينولات}، يحابيه التسخين (\cref{prop:b2:enolates-aldol:aldol-equilibrium}).
\textbf{18.} كيتون مشبع وكيتون $\alpha,\beta$-غير مشبع (هكسينون حلقي).
\textbf{19.} كربون التحام الحلقتين الحامل لمجموعة الميثيل: له أربعة بدائل مختلفة.
\textbf{20.} لا: فالديون لاكيرالي، ويُهاجَم كربونيلاه بالتساوي (فهما صورتان مرآويتان أحدهما للآخر
بالنسبة إلى السلسلة)؛ ومع كواشف لاكيرالية يكون الناتج راسيميًا. والحفازات الأمينية الكيرالية
تحابي أحد الكربونيلين وتعطي أحد المتماكبين المرآويين بفائض (مجلد السنة الثالثة).
\textbf{21.} الديون \ce{C7H10O2}: \qty{126.155}{g/mol}؛ \ce{C4H6O}: \qty{70.091}{g/mol}؛ \ce{C11H14O2}:
\qty{178.231}{g/mol}.
\textbf{22.} \ce{C7H10O2 + C4H6O -> C11H14O2 + H2O}: C 11، وH 16، وO 3 في كل طرف.
\textbf{23.} $10.0/126.155 = \qty{0.07927}{mol} = \qty{79.3}{mmol}$.
\textbf{24.} $0.07927 \times 178.231 \times 0.70 = \boldsymbol{\approx \qty{9.89}{g}}$ من
كيتون فيلاند--ميشر.
\end{solution}
